\documentclass[a4paper,10pt]{article}
\usepackage[utf8]{inputenc}
\usepackage[T1]{fontenc}
\usepackage{lmodern}
\usepackage[margin=1.5cm]{geometry}
\usepackage{titlesec}
\usepackage{enumitem}
\usepackage{hyperref}
\usepackage{xcolor}
\usepackage{fontawesome5}
\usepackage{tabularx}

% Colors
\definecolor{primary}{HTML}{000000}
\definecolor{accent}{HTML}{2E5A88}
\definecolor{link}{HTML}{1565C0}

\hypersetup{
    colorlinks=true,
    linkcolor=link,
    urlcolor=link,
    pdfauthor={Gabriele Vianello},
    pdftitle={CV Gabriele Vianello}
}

\titleformat{\section}{\Large\bfseries\color{primary}}{}{0em}{}[\vspace{-0.5em}\rule{\textwidth}{0.4pt}]
\titlespacing*{\section}{0pt}{1em}{0.5em}

\newcommand{\cventry}[4]{
  \noindent\textbf{#1} \hfill \textit{#2}\\
  \textit{#3} \hfill #4\vspace{0.3em}
}

\newcommand{\projectentry}[3]{
  \noindent\textbf{#1} | \textit{#2}\\
  #3\vspace{0.3em}
}

\pagestyle{empty}

\begin{document}

\begin{center}
  {\Huge\bfseries Gabriele Vianello} \\[0.5em]
  {\Large\color{accent} AI \& Data Science Engineer} \\[0.5em]
  \faEnvelope~\href{mailto:vianello.tech@gmail.com}{vianello.tech@gmail.com} \quad
  \faGithub~\href{https://github.com/Vinello28}{Vinello28} \quad
  \faLinkedin~\href{https://www.linkedin.com/in/gabriele-vianello-476a331a1}{LinkedIn} \quad
  \faMapMarker~Italia
\end{center}

\vspace{0.5em}

\section{Profilo}
Studente magistrale in Ingegneria Informatica (AI \& Data Science). 
Appassionato di Machine Learning, architetture Zero-Trust e Data Engineering. 
Esperienza pratica nella progettazione di pipeline ML, fine-tuning di modelli (LLM/VLM) e sviluppo di backend ad alte prestazioni.

\section{Competenze Tecniche}
\begin{tabularx}{\textwidth}{@{}l X@{}}
  \textbf{AI \& Data Science:} & PyTorch, TensorFlow, Hugging Face, Unsloth, Scikit-learn, Pandas, PyTorch Geometric, QLoRA \\
  \textbf{Linguaggi:} & Python, Go, C\#, TypeScript, C, SQL \\
  \textbf{Serving \& DevOps:} & Docker, FastAPI, vLLM, ONNX Runtime, Ollama, MongoDB, PostgreSQL, DuckDB, GitHub Actions \\
  \textbf{Frontend:} & React, Vite, TailwindCSS, HTMX \\
\end{tabularx}

\section{Progetti Selezionati}

\projectentry{Graphagate}{AI Engineer \& ML Researcher}{Temporal Graph Network per il rilevamento di anomalie non supervisionato su flussi di accesso Zero-Trust. Addestrato con negative sampling per classificare eventi in tempo reale (anomalie contestuali, policy, lateral-movement). Servito tramite microservizio FastAPI. \textit{Tecnologie: PyTorch, PyTorch Geometric, FastAPI, Docker}}

\projectentry{ZTALeaks}{Security \& AI Engineer}{Implementazione a microservizi di un'architettura Zero Trust (NIST SP 800-207). Separazione PEP/PDP su segmenti di rete isolati, con decisioni di accesso basate su policy OPA Rego e risk score AI in real-time. \textit{Tecnologie: Go, OPA, Envoy, MongoDB, Docker, Snort}}

\projectentry{synThor}{Data Engineer \& ML Engineer}{Generatore di dataset sintetici per il fine-tuning di modelli Vision-Language su estrazione documentale. Produce coppie allineate immagine/JSON per 4 tipi di documenti in 11 lingue, mantenendo fisso lo schema di estrazione. \textit{Tecnologie: Python, PyMuPDF, ReportLab, Unsloth}}

\projectentry{Pack-a-Punch}{AI Engineer}{Classificatore binario (AI-generated vs human-written test) su base bert-base-italian-xxl-cased. Inferenza ottimizzata tramite ONNX Runtime e accelerazione CUDA/Metal ($\sim$110 req/s), con knowledge distillation. Servizio FastAPI containerizzato. \textit{Tecnologie: Python, PyTorch, ONNX Runtime, FastAPI, Docker}}

\projectentry{Charge-a-Sloth}{AI Engineer}{Pipeline ottimizzata per il fine-tuning di modelli vision-language con Unsloth e QLoRA. Addestramento PEFT 4-bit su dataset parquet locali convertiti in chat templates. \textit{Tecnologie: PyTorch, Unsloth, QLoRA, Hugging Face, CUDA}}

\projectentry{Pack-a-Mail}{Backend Engineer}{Servizio per l'invio massivo di email ad alto throughput, scritto in Go. Architettura Clean (ports \& adapters), backend SQLite CGO-free, dashboard HTMX integrata e hashing Argon2id. Distribuito come singolo eseguibile. \textit{Tecnologie: Go, SQLite, HTMX, Clean Architecture}}

\section{Formazione}
\cventry{Laurea Magistrale in Ingegneria Informatica (AI \& Data Science)}{In corso}{Università Politecnica delle Marche}{Ancona}
\vspace{-0.5em}
\begin{itemize}[leftmargin=1.5em, itemsep=0pt, topsep=0pt]
  \item Corsi chiave: Deep Learning, Computer Vision, NLP, Data Mining, Network Science.
\end{itemize}

\vspace{0.5em}
\cventry{Laurea Triennale in Ingegneria dei Sistemi Informatici}{Completata}{Politecnico di Milano}{Cremona/Milano}
\vspace{-0.5em}
\begin{itemize}[leftmargin=1.5em, itemsep=0pt, topsep=0pt]
  \item Solida formazione in fondamenti dell'informatica, algoritmi e architetture software.
\end{itemize}

\vspace{0.5em}
\cventry{Diploma di Perito Informatico}{Completato}{Istituto Tecnico Industriale V. Volterra}{Ancona}
\vspace{-0.5em}
\begin{itemize}[leftmargin=1.5em, itemsep=0pt, topsep=0pt]
  \item Indirizzo Informatica e Telecomunicazioni, specializzazione Informatica.
\end{itemize}

\end{document}
